%%
%%Analyse en ontwerp van de oplossing
%%

\chapter{Analyse en ontwerp van de oplossing}
\label{ch:analyseontwerp}

\section{Doelstellingen van de oplossing}

Op basis van de analyse van het huidige onboarding- en offboardingsprocedures werd vastgesteld dat verschillende stappen binnen de ICT-dienst repetitief en manueel uitgevoerd worden. Hierdoor neemt het verwerken van nieuwe medewerkers relatief veel tijd in beslag en ontstaat een verhoogde kans op menselijke fouten. Daarnaast zijn verschillende diensten afhankelijk van elkaar voor het uitvoeren van bepaalde configuraties, waardoor vertragingen kunnen ontstaan wanneer terugkoppelingen ontbreken.

Het doel van dit onderzoek is daarom het ontwerpen van een oplossing die een deel van het onboarding- en van het offboardingsproces automatiseert binnen de bestaande ICT-omgeving van het gemeentebestuur van Berlare. De focus ligt hierbij voornamelijk op het automatisch verwerken van gebruikersgegevens en het automatisch aanmaken of uitschakelen van gebruikersaccounts binnen de Active Directory-omgeving.

Daarnaast moet de oplossing zorgen voor een betere opvolging van uitgevoerde taken en een efficiëntere samenwerking tussen verschillende diensten binnen het onboarding- en offboardingsproces. Het uiteindelijke doel is om de administratieve belasting te verminderen, de foutmarge te verkleinen en de doorlooptijd van beide processen te beperken.

\section{Requirements van de oplossing}

\subsection{Functionele requirements}

\subsubsection{Must-have requirements}

De belangrijkste doelstelling van deze bachelorproef is het ontwikkelen van een geautomatiseerd systeem waarbij gebruikersaccounts automatisch aangemaakt worden op basis van gegevens die aangeleverd worden door de personeelsdienst.

De focus ligt hierbij voornamelijk op het automatisch aanmaken van gebruikersaccounts binnen Active Directory, het instellen van relevante attributen binnen de lokale infrastructuur, het toekennen van Microsoft 365-licenties en het automatisch voorzien van basisrechten afhankelijk van de dienst van de medewerker.

Daarnaast moet de oplossing voorzien in een vorm van terugkoppeling zodat de status van uitgevoerde onboardings beschikbaar is voor andere betrokken diensten binnen de organisatie.

\subsubsection{Should-have requirements}

Naast de basisfunctionaliteiten zijn er verschillende uitbreidingen die bij voorkeur aanwezig zijn binnen de oplossing. Hieronder vallen onder andere het automatisch toekennen van rechten op gedeelde postvakken en het automatisch toevoegen van gebruikers aan distributielijsten binnen Microsoft 365.

Deze functionaliteiten verhogen de gebruiksvriendelijkheid van het systeem en verminderen bijkomende manuele configuraties binnen de ICT-dienst.

\subsubsection{Could-have requirements}

Binnen deze categorie vallen functionaliteiten die een meerwaarde bieden, maar niet noodzakelijk zijn voor de basiswerking van de oplossing.

Een eerste uitbreiding is het automatisch versturen van een welkomstmail naar nieuwe medewerkers met praktische informatie en handleidingen. Daarnaast behoort ook het automatiseren van het offboardingproces tot deze categorie.

Hoewel offboarding een belangrijk onderdeel vormt van identiteits- en toegangsbeheer binnen organisaties, lag de primaire focus van deze bachelorproef op onboarding. Indien mogelijk werd echter ook een gedeeltelijke automatisering van offboarding opgenomen binnen de oplossing.

\subsubsection{Would-have requirements}

De laatste categorie omvat uitbreidingen die wenselijk zijn, maar momenteel moeilijk realiseerbaar zijn binnen de bestaande infrastructuur.

Hieronder valt het volledig automatiseren van rechtenbeheer op de file server. Binnen de huidige omgeving bestaat een complexe rechtenstructuur waarbij bepaalde toegangen historisch rechtstreeks aan individuele gebruikers toegewezen werden in plaats van via veiligheidsgroepen.

Door deze inconsistente structuur is het momenteel niet mogelijk om alle rechten volledig automatisch en betrouwbaar toe te wijzen. Een volledige herstructurering van de bestaande rechtenomgeving zou hiervoor noodzakelijk zijn.

\subsection{Niet-functionele requirements}

De oplossing moet compatibel zijn met de bestaande hybride infrastructuur van het gemeentebestuur van Berlare. Zowel lokale on-premisesystemen als cloudtoepassingen binnen Microsoft 365 moeten ondersteund worden.

Daarnaast moet de oplossing veilig functioneren binnen de bestaande infrastructuur. Enkel bevoegde medewerkers mogen toegang hebben tot de automatiseringsflows en gevoelige configuraties.

Verder moet de oplossing voldoende flexibel en onderhoudbaar zijn zodat toekomstige uitbreidingen en aanpassingen mogelijk blijven.

\section{Analyse van automatiseerbare processen}

\subsection{Automatiseerbare taken}

Binnen het huidige onboardingproces werden verschillende taken geïdentificeerd die zich lenen tot automatisering. Het gaat hierbij voornamelijk om digitale configuraties die volgens vaste patronen uitgevoerd worden binnen de ICT-omgeving.

Een eerste belangrijke taak is het automatisch aanmaken van gebruikersaccounts binnen Active Directory. Momenteel gebeurt dit handmatig waarbij gebruikersgegevens individueel ingevoerd worden. Daarnaast worden gebruikers eveneens manueel toegevoegd aan beveiligingsgroepen en distributielijsten.

Ook het toekennen van licenties binnen Entra ID gebeurt momenteel afzonderlijk per gebruiker. Omdat deze toewijzingen grotendeels gebaseerd zijn op functie en dienst, kan dit proces geautomatiseerd worden door gebruikers automatisch aan de correcte groepen toe te voegen.

Verder werden ook terugkerende communicatieve taken geïdentificeerd. Zo kunnen meldingen of terugkoppelingen automatisch verstuurd worden wanneer bepaalde configuraties succesvol uitgevoerd werden. Hierdoor kunnen andere diensten sneller verder werken met hun eigen configuraties.

Naast onboarding werden ook verschillende onderdelen van het offboardingproces geïdentificeerd die zich lenen tot automatisering. Wanneer een medewerker de organisatie verlaat, moeten verschillende toegangen en configuraties aangepast of verwijderd worden binnen zowel Active Directory als Microsoft 365.

Momenteel gebeurt dit grotendeels manueel door de ICT-dienst. Hierbij worden gebruikersaccounts gedeactiveerd, groepslidmaatschappen verwijderd, mailboxrechten aangepast en licenties ingetrokken. Omdat deze stappen volgens vaste procedures verlopen, kunnen ook deze processen gedeeltelijk geautomatiseerd worden.

Daarnaast is offboarding vanuit beveiligingsoogpunt een belangrijk proces. Het tijdig verwijderen van toegangsrechten vermindert het risico op ongewenste toegang tot interne systemen en gegevens nadat een medewerker uit dienst gegaan is.


\subsection{Niet-automatiseerbare taken}

Niet alle onderdelen van het onboardingproces lenen zich tot automatisering. Bepaalde taken vereisen fysieke tussenkomst of individuele beoordeling door medewerkers van de ICT-dienst.

Zo gebeurt het selecteren van geschikte hardware momenteel op basis van de functie en de specifieke noden van de medewerker. Daarnaast vereisen ook de installatie van toestellen, het aansluiten van randapparatuur en bepaalde softwareconfiguraties manuele handelingen.

Ook het fysiek verwerken van documenten, zoals het afdrukken en overhandigen van tijdelijke wachtwoorden, valt buiten de scope van dit onderzoek. Hiervoor zouden bijkomende gespecialiseerde systemen noodzakelijk zijn.

De focus van dit onderzoek ligt daarom voornamelijk op het automatiseren van digitale configuraties en identiteitsbeheer binnen de bestaande IT-omgeving.

Ook binnen het offboardingproces blijven bepaalde taken afhankelijk van fysieke tussenkomst door de ICT-dienst. Zo kunnen het ophalen van hardware, het fysiek controleren van toestellen en het volledig herinstalleren van laptops niet automatisch uitgevoerd worden binnen de scope van deze oplossing.


\section{Technische vereisten}

Binnen de bestaande infrastructuur van het gemeentebestuur wordt gebruikgemaakt van een hybride omgeving waarbij zowel cloudtoepassingen als lokale on-premisesystemen aanwezig zijn. Hierdoor moet de oplossing kunnen communiceren met beide omgevingen.

De gebruikersaccounts worden lokaal beheerd binnen Active Directory, terwijl bepaalde cloudtoepassingen gebruikmaken van Entra ID. Hierdoor moet de oplossing rekening houden met synchronisatie tussen beide platformen.

Daarnaast vormt compatibiliteit met bestaande Microsoft-diensten een belangrijke vereiste binnen het project. Binnen de organisatie wordt reeds gebruikgemaakt van Microsoft Power Automate, Microsoft 365 en PowerShell-gebaseerde beheertools. Hierdoor moet de oplossing integreren met deze bestaande infrastructuur.

Verder moet de oplossing voldoende flexibel zijn om rekening te houden met verschillende diensten, sites en organisatorische structuren binnen de Active Directory-omgeving.

Daarnaast moet de oplossing voldoende veilig functioneren binnen de bestaande infrastructuur. Omdat gevoelige configuraties uitgevoerd worden binnen zowel Active Directory als Microsoft 365, moet toegangsbeheer correct afgeschermd worden en mogen enkel bevoegde medewerkers de automatiseringsflows kunnen uitvoeren.

\section{Keuze van technologieën}

Binnen het onderzoek werden verschillende automatiseringsplatformen onderzocht die gebruikt kunnen worden voor workflowautomatisering en systeemintegratie.

Platformen zoals Zapier, Make en Activepieces bieden mogelijkheden voor cloudgebaseerde automatisering en integratie tussen verschillende toepassingen. Deze platformen richten zich voornamelijk op algemene workflowautomatisering tussen cloudservices.

Binnen dit project vormde ondersteuning voor lokale PowerShell-uitvoering echter een belangrijke vereiste. Omdat gebruikersaccounts lokaal binnen Active Directory aangemaakt moeten worden, moet de oplossing kunnen communiceren met on-premisesystemen.

Daarnaast werd binnen de organisatie reeds gebruikgemaakt van infrastructuur gebaseerd op Microsoft en bestaande automatiseringsprojecten binnen Power Automate. Hierdoor bood Power Automate belangrijke voordelen op vlak van compatibiliteit en integratie met bestaande Microsoft-diensten.

Verder ondersteunt Power Automate de combinatie van cloudflows en lokale desktopflows via Power Automate Desktop en Machine Runtime, waardoor communicatie met lokale systemen mogelijk wordt binnen een hybride omgeving.

Op basis van deze vereisten werd gekozen om de verdere implementatie uit te werken met Microsoft Power Automate in combinatie met PowerShell-gebaseerde automatisering.

Daarnaast moet de oplossing niet alleen nieuwe accounts kunnen aanmaken, maar ook bestaande gebruikersaccounts veilig kunnen aanpassen of deactiveren tijdens het offboardingproces. Hiervoor moet eveneens communicatie mogelijk zijn met Exchange Online voor het beheren van mailboxrechten, distributiegroepen en gedeelde postvakken.

\section{Ontwerp van de oplossing}
\subsection{Algemene oplossing}

De ontworpen oplossing bestaat uit verschillende componenten die samenwerken binnen zowel de cloudomgeving als de lokale infrastructuur van de gemeente.

De invoer van gebruikersgegevens gebeurt via Microsoft Forms. Wanneer een nieuwe aanvraag wordt verzonden, worden de gegevens verwerkt binnen een cloudflow in Power Automate. Deze flow verwerkt de ontvangen gegevens en bereidt deze voor op verdere verwerking binnen de lokale omgeving.

Omdat Active Directory lokaal beheerd wordt, moet de cloudomgeving kunnen communiceren met een on-premisesysteem. Hiervoor wordt gebruikgemaakt van Power Automate Desktop en Power Automate Machine Runtime. Deze componenten vormen de brug tussen de cloudflow en de lokale serveromgeving.

Binnen de lokale omgeving wordt een PowerShell-script uitgevoerd dat verantwoordelijk is voor het aanmaken of uitschakelen van gebruikersaccounts, het toewijzen of verwijderen van groepen en het configureren van verschillende Active Directory-attributen.

Bij onboarding moet eerst gewacht worden tot de synchronisatie tussen Active Directory en Entra ID voltooid is voordat cloudconfiguraties uitgevoerd kunnen worden. Bij offboarding is deze synchronisatie minder kritisch omdat het uitschakelen van accounts voornamelijk lokaal verwerkt wordt en de blokkering van cloudtoegang automatisch volgt via synchronisatie.

Nadien worden de distributielijsten en gedeelde postvakken toegevoegd of verwijderd bij de gebruiker. ODe configuratie van distributielijsten en gedeelde postvakken gebeurt via PowerShell-scripts die communiceren met Exchange Online. Hiervoor werd gekozen omdat Power Automate momenteel geen volwaardige connectoren aanbiedt voor het beheren van mailboxrechten binnen Exchange Online.

Als laatste stap binnen het configuratieproces wordt automatisch een welkomstmail verzonden naar de nieuwe medewerker. Deze mail bevat verschillende handleidingen en praktische informatie die noodzakelijk zijn om vlot van start te kunnen gaan tijdens de eerste werkdag. Hierbij kan gedacht worden aan informatie over het verbinden met het draadloze netwerk, het gebruik van printers en andere interne ICT-toepassingen.

Na het voltooien van alle configuratiestappen wordt een terugkoppeling uitgevoerd binnen het opvolgsysteem. Omdat het algemene opvolgingssysteem van een parallel project momenteel nog in ontwikkeling is, werd tijdelijk gekozen voor een tussenoplossing op basis van Excel. Hierbij wordt de status van de verwerking automatisch aangepast binnen hetzelfde Excelbestand dat tijdens de testfase gebruikt werd. Op deze manier kan eenvoudig gecontroleerd worden of de volledige flow succesvol werd uitgevoerd en de gebruiker correct werd aangemaakt.

Figuur~\ref{fig:ArchitectuurOnboarding} toont een vereenvoudigd overzicht van de algemene architectuur van de oplossing voor onboarding.

\begin{figure}[htbp]
    \centering
    \includegraphics[width=1\textwidth]{Architectuur.png}
    \caption[Power Automate flow]{Vereenvoudigde voorstelling van de Power Automate Flow voor onboarding}
    \label{fig:ArchitectuurOnboarding}
\end{figure}

\subsection{Ontwerp van de onboardingsflow}

De automatiseringsflow werd ontworpen op basis van de bestaande onboardingprocedure binnen de ICT-dienst. Hierbij werd een verwerking voorzien waarbij gebruikersgegevens automatisch vertaald worden naar configuraties binnen Active Directory en Entra ID.

Binnen Microsoft Forms worden verschillende gegevens verzameld, waaronder voornaam, achternaam, functietitel, dienst, site, telefoonnummer en contractinformatie. Deze gegevens worden vervolgens verwerkt binnen Power Automate.

Omdat bepaalde vragen binnen het formulier afhankelijk zijn van eerder gekozen antwoorden, moet de flow rekening houden met dynamische gegevensstructuren. Hiervoor worden conditionele controles gebruikt om de correcte diensten en bijhorende OU-structuren te bepalen.

Daarnaast moet de oplossing rekening houden met verschillen tussen cloud- en on-premisesystemen. Hierdoor werd gekozen voor een architectuur waarbij de cloudflow verantwoordelijk is voor gegevensverwerking en orchestration, terwijl lokale PowerShell-scripts instaan voor de configuratie van Active Directory.

\subsection{Ontwerp van de offboardingflow}

Naast het automatiseren van onboarding werd eveneens een geautomatiseerde offboardingflow ontworpen. Deze flow heeft als doel om gebruikersaccounts en toegangsrechten automatisch af te bouwen wanneer een medewerker uit dienst gaat.

De offboarding start vanuit een afzonderlijk Microsoft Forms-formulier waarin de personeelsdienst de gegevens van de medewerker invoert, waaronder het e-mailadres, de einddatum van het contract en de reden van uitdiensttreding.

Omdat gebruikers pas gedeactiveerd mogen worden nadat de laatste werkdag verstreken is, werd een wachmechanisme voorzien binnen de cloudflow van Power Automate. De flow blijft actief totdat de opgegeven einddatum bereikt werd.

Na het bereiken van deze datum wordt automatisch een desktopflow gestart die verschillende PowerShell-scripts uitvoert binnen zowel Exchange Online als Active Directory. Hierbij wordt de mailbox van de gebruiker omgezet naar een gedeeld postvak, worden mailboxrechten verwijderd, groepslidmaatschappen aangepast en wordt het Active Directory-account gedeactiveerd.

Daarnaast worden Microsoft 365-groepen en licenties automatisch verwijderd zodat ongebruikte cloudlicenties opnieuw beschikbaar worden binnen de organisatie.

Door deze automatisering wordt het offboardingproces consistenter uitgevoerd en wordt het risico verminderd dat voormalige medewerkers nog toegang behouden tot interne systemen of gegevens.

\subsection{OU-structuur binnen Active Directory}

Binnen de Active Directory-omgeving van de gemeente wordt gebruikgemaakt van een hiërarchische structuur van Organizational Units (OU's). Gebruikers worden hierbij ondergebracht op basis van site en dienst.

Afhankelijk van de organisatorische structuur kunnen diensten zich op verschillende niveaus binnen de directorystructuur bevinden. Sommige sites bevatten afzonderlijke OU's per dienst, terwijl andere diensten gezamenlijk ondergebracht worden binnen één gedeelde OU.

Zo bevat de site ``Gemeentehuis'' afzonderlijke OU's voor diensten zoals personeelsdienst, vergunningen en informatica. Andere locaties, zoals de onthaalpoort aan het Donkmeer, groeperen meerdere diensten binnen dezelfde OU omdat deze organisatorisch nauwer samenwerken.

Door deze variatie binnen de directorystructuur moet de automatiseringsoplossing rekening houden met verschillende mappings tussen diensten en hun overeenkomstige OU-locaties.

Figuur~\ref{fig:OrganizationalUnit} toont een vereenvoudigd overzicht van de gebruikte OU-structuur binnen Active Directory.

\clearpage

\begin{figure}[htbp]
    \centering
    \includegraphics[width=0.4\textwidth]{OrganizationalUnit.png}
    \caption[OU structuur]{Vereenvoudigde voorstelling van de OU-structuur binnen de Active Directory-omgeving van de gemeente.}
    \label{fig:OrganizationalUnit}
\end{figure}

\subsection{Synchronisatie tussen Active Directory en Entra ID}

Binnen de hybride omgeving van de gemeente worden gebruikersaccounts eerst lokaal aangemaakt binnen Active Directory. Vervolgens worden deze accounts automatisch gesynchroniseerd naar Entra ID.

Deze synchronisatie gebeurt echter niet onmiddellijk maar volgens vaste synchronisatiecycli. Hierdoor ontstaat een vertraging tussen het aanmaken van een gebruiker binnen Active Directory en de beschikbaarheid van deze gebruiker binnen Entra ID.

Omdat bepaalde cloudconfiguraties, zoals licentietoewijzingen, pas uitgevoerd kunnen worden nadat de gebruiker beschikbaar is binnen Entra ID, moet de automatiseringsflow rekening houden met deze wachttijd.

Hiervoor werd een controlemechanisme ontworpen waarbij periodiek gecontroleerd wordt of de gebruiker reeds beschikbaar is binnen Entra ID. Pas wanneer de synchronisatie succesvol voltooid werd, kunnen bijkomende cloudconfiguraties uitgevoerd worden.

Ook tijdens offboarding blijft deze synchronisatie belangrijk zodat wijzigingen aan gebruikersaccounts en groepslidmaatschappen correct doorstromen naar de cloudomgeving.

\section{Testomgeving en validatieaanpak}

Tijdens de eerste ontwikkelingsfase was nog geen volledige productieomgeving beschikbaar voor het uitvoeren van geautomatiseerde cloudflows. Daarnaast vereisen bepaalde onderdelen van Power Automate bijkomende licenties.

Om de ontwikkeling toch te kunnen opstarten, werd eerst gewerkt met een testomgeving waarbij gebruikersgegevens vanuit een Excelbestand verwerkt werden. De kolommen binnen dit bestand werden opgesteld volgens dezelfde structuur als de uiteindelijke variabelen binnen Power Automate.

Deze aanpak maakte het mogelijk om de PowerShell-scripts afzonderlijk te ontwikkelen en te testen zonder afhankelijk te zijn van de volledige cloudomgeving. Daarnaast kon hierdoor gecontroleerd worden of gebruikers correct aangemaakt werden en de juiste groepslidmaatschappen ontvingen.

De tijdelijke testomgeving vormde hierdoor een belangrijke tussenfase binnen het ontwikkelingsproces van de uiteindelijke oplossing.



